\documentclass[12pt,a4paper]{report}

\usepackage[utf8]{inputenc}
\usepackage[T1]{fontenc}
\usepackage[spanish]{babel}
\usepackage{listings}
% \usepackage{xcolor}
% \usepackage{geometry}
\usepackage{hyperref}
\usepackage{parskip}
\usepackage{titlesec}
\usepackage{booktabs}
\usepackage{tcolorbox}
\usepackage{amsmath}
\usepackage{array}

% \geometry{margin=2.5cm}
\usepackage[dvipsnames]{xcolor}
\usepackage[lmargin=2.5cm, rmargin=5cm]{geometry}
\input{word-comments.tex}



\definecolor{lisp-bg}{RGB}{248,248,252}
\definecolor{lisp-kw}{RGB}{0,0,180}
\definecolor{lisp-str}{RGB}{160,0,0}
\definecolor{box-blue}{RGB}{240,248,255}
\definecolor{box-blue-border}{RGB}{100,149,237}
\definecolor{box-green}{RGB}{240,255,240}
\definecolor{box-green-border}{RGB}{60,150,60}
\definecolor{box-yellow}{RGB}{255,253,230}
\definecolor{box-yellow-border}{RGB}{180,150,0}

\lstdefinelanguage{DSL}{
  keywords={def-table, libro, hoja, tabla, load, defparameter,
            conditional-rendering, exists, nav, ultima-fila, primera-fila,
            xl-generate, xl-run-generated},
  morekeywords=[2]{countif, counta, sum-range, str, col, trange, param,
                   add, divide, promedio, gt, _if},
  keywordstyle=\color{lisp-kw}\bfseries,
  keywordstyle=[2]\color{box-green-border},
  stringstyle=\color{lisp-str},
  sensitive=false,
  basicstyle=\ttfamily\small,
  backgroundcolor=\color{lisp-bg},
  frame=single,
  framerule=0.4pt,
  rulecolor=\color{gray!40},
  breaklines=true,
  numbers=left,
  numberstyle=\tiny\color{gray},
  numbersep=8pt,
  xleftmargin=20pt,
  aboveskip=6pt,
  belowskip=6pt,
  showstringspaces=false,
  tabsize=2,
}

\lstset{language=DSL}

\tcbuselibrary{skins,breakable}

\newtcolorbox{infobox}[1]{
  colback=box-blue, colframe=box-blue-border,
  fonttitle=\bfseries, title=#1,
  breakable, sharp corners, left=6pt, right=6pt, top=4pt, bottom=4pt
}

\newtcolorbox{advertencia}[1]{
  colback=box-yellow, colframe=box-yellow-border,
  fonttitle=\bfseries, title=#1,
  breakable, sharp corners, left=6pt, right=6pt, top=4pt, bottom=4pt
}

\titleformat{\chapter}[hang]{\large\bfseries}{\thechapter.}{0.5em}{}[\vspace{2pt}\hrule]
\titleformat{\section}[hang]{\normalsize\bfseries}{\thesection}{0.5em}{}
\titleformat{\subsection}[hang]{\normalsize\itshape\bfseries}{\thesubsection}{0.5em}{}

\hypersetup{hidelinks}
\setlength{\parskip}{6pt}

\begin{document}

\chapter{Tutorial básico: tablas, fórmulas y resaltado}
\label{cap:tutorial-basico}

Este capítulo presenta los elementos fundamentales del DSL mediante un
ejemplo mínimo.  Al terminar sabrá definir una tabla con columnas, añadir
una columna calculada, condicionar el resaltado de celdas y generar el
archivo de salida.

\section{Qué vamos a construir}

Queremos una tabla con cuatro columnas:

\notaparaelautor{le voy a poner los separadores de las columnas a la tabla.}

\notaparaelautor{También voy a resaltar en rojo los que tengan más de 80}

\begin{center}
\begin{tabular}{|l|c|c|r|}
\toprule
\textbf{Nombre} & \textbf{Nota 1} & \textbf{Nota 2} & \textbf{Promedio} \\
\midrule
\textcolor{red}{Ana}    & 85 & 90 & 87.5 \\
Luis   & 70 & 75 & 72.5 \\
\textcolor{red}{Elena}  & 92 & 88 & 90.0 \\
Carlos & 60 & 65 & 62.5 \\
\textcolor{red}{Sofía}  & 78 & 82 & 80.0 \\
\textcolor{red}{Pedro}  & 95 & 91 & 93.0 \\
\bottomrule
\end{tabular}
\end{center}

Además:
\begin{itemize}
  \item La columna \textbf{Promedio} no se escribe a mano: se calcula
    automáticamente como el promedio de las dos notas.
  \item Hay un \textbf{número fijo} (80, por ejemplo) almacenado en el
    libro.
  \item Los nombres de las personas cuyo promedio supera ese número
    aparecen \textbf{resaltados en rojo}.
\end{itemize}

\section{Estructura del programa}

El programa completo consta de tres partes, que escribiremos \comment{en un solo
archivo \texttt{tutorial-basico.lisp}}{¿que en qué carpeta lo creo?}:

\begin{enumerate}
  \item Los \textbf{datos} de entrada (nombres y notas).
  \item La \textbf{definición de la tabla} con su columna calculada y su
    regla de resaltado.
  \item El \textbf{libro} que reúne todo y las instrucciones para generar
    la salida.
\end{enumerate}

\notaparaelautor{Mmmm... aquí sugeriría hacerlo al revés... empezar por la tabla primero y después poner los datos... ¿no? La alternativa a eso es que expliques entonces por qué los datos de una tabla están en una lista de common lisp... Recuerda que el tutorial es para alguien que no tenga idea de nada de nada...}

\notaparaelautor{Deja ver cómo sigue la cosa, pero yo pensaría que la idea es explicar primero la estructura de la tabla (con el def-table, que supuestamemente crea una función... ¿no?) y después cuando vayamos a llamar esa función, si no le pasamos datos, crearía una tabla sin datos, que no está mal, lo llenas a mano y ya, pero es más \textit{tocào} cuando le pones datos... y ahí presentas los datos en forma de lista.}

\notaparaelautor{En cualquier caso, hay un problema más grande, que es que no me has dicho qué ficheros crear y desde dónde correr sbcl... por aquí falta una instrucción que sea crea el fichero en la misma carpeta donde están todas las cosas, y entonces... y sigues por ahí...}

\section{Los datos}

Los datos se guardan en una variable de Lisp con
\texttt{defparameter}.  Cada fila es una lista con cuatro
\comment{celdas}{será elementos, ¿no? ¿o qué significa que una lista
  tenga celdas?} (una por columna).  Las \comment{\queesesto{primeras tres filas son de
prefijo}}{Antes de esto, en algún momento deberías explicarme qué son los prefijos esos}: dos filas en blanco y una fila de encabezados.  Las filas
siguientes son los datos propiamente dichos.

\begin{lstlisting}[caption={tutorial-basico.lisp --- datos de entrada.}]
(defparameter *DATOS*
  '(("" "" "" "")
    ("" "" "" "")
    ("Nombre" "Nota 1" "Nota 2" "Promedio")
    ("Ana"     "85" "90" "")
    ("Luis"    "70" "75" "")
    ("Elena"   "92" "88" "")
    ("Carlos"  "60" "65" "")
    ("Sofia"   "78" "82" "")
    ("Pedro"   "95" "91" "")))
\end{lstlisting}

Observe que la columna \texttt{Promedio} se \comment{deja vacía (\texttt{""})}{Por algún motivo aquí no salen las comillas :-/... y lo que veo es \textit{una lista vacía} en Common Lisp, que me confunde un poquito...} en
los datos, porque su valor lo calculará \comment{el DSL automáticamente}{¿el DSL o el excel cuando se cree después? Aquí sería bueno separar el DSL del producto final...}.

\begin{advertencia}{Formato de las celdas vacías}
  Las celdas que corresponden a columnas calculadas se marcan con
  \comment{\texttt{""}}{Decididamente esta \textit{cadena vacía} no se ve. Hay que buscar una forma que de que vean.} (cadena vacía).  \queesesto{El generador} las sobreescribe \queesesto{con la
  fórmula correspondiente}.
\end{advertencia}

\agregaesto{TRANSICIÓN A LO QUE SIGUE.}

\section{Definir la tabla}

La macro \texttt{def-table} declara \queesesto{una plantilla de
  tabla}.  Le damos un nombre y la lista de columnas.  Cada columna se
escribe como un par \texttt{(nombre "Etiqueta")}, donde
\texttt{nombre} es el \queesesto{identificador interno} (un símbolo del DSL) y
\texttt{"Etiqueta"} es el texto que se mostrará \comment{como encabezado}{¿cómo encabezado de qué?}.

\agregaesto{Aquí sería bueno poner un ejemplito de eso de los encabezados y los identificadores internos. Cuando me refiero a un ejemplito, no me refiero a la tabla que viene ahí a continuación sino a algo explicado con palabras... para entender la idea, y después poderlo \textit{entender} en el ejemplo.}

\begin{lstlisting}[caption={Definición de la tabla.}]
(def-table tabla-notas
  ((nombre "") (nota1 "") (nota2 "") (promedio ""))
  :cell-height 1
  :cell-width  1)
\end{lstlisting}

\notaparaelautor{Yo sé que en el tutorial no me has mandado a hacerlo, pero en common lisp, la idea es ir creando las cosas y probándolas inmediatamente... si trato de ejecutar ese fragmento de código, me da este error}

\begin{verbatim}
The variable TABLA-NOTAS is unbound.
   [Condition of type UNBOUND-VARIABLE]

\end{verbatim}

Por ahora la tabla solo declara las columnas, sin fórmulas ni resaltados. \notaparaelautor{Aquí yo sugeriría decir que si creamos un libro con esta tabla solo estarían los encabezados... y ya, más nada.}
A continuación añadiremos esas características.

\notaparaelautor{Mmm.... a lo mejor conviene que el tutorial tenga tres partes: una primera en la que solo se crean los encabezados de la tabla, ni siquiera con datos. y con eso, crear un libro y exportarlo y todo. Incluso, mostrar el resultado del excel, poniendo unas foticos. Ese sería el tutorial ultra básico para confirmar que el sistema funciona. Después, vamos a ponerle datos, y después las otras dos cosas.}

\notaparaelautor{Para eso antes de tirar una línea de código debes explicar lo que se puede hacer con esto: crear tablas, que se ponen en hojas, que a su vez, se ponen en libros. Todo eso se hace en lisp. Después, esos \textit{libros} se exportan a alguno de los backend disponibles, y después, puedes experimentar con esas tablas, hojas y libros, en la arquitectura en la que se hayan creado.}

\section{Columna calculada}

Queremos \comment{que \texttt{promedio}}{¿qué es \texttt{promedio}? ¿o te refieres a los elementos de la columna promedio?} \comment{se calcule solo}{¿qué significa esto?}: \[
\text{promedio} = \frac{\text{nota1} + \text{nota2}}{2}
\]

\notaparaelautor{¿qué son nota1 y nota2?}

En lenguaje natural diríamos: ``\comment{para cada fila}{¿qué significa esto?}, promedia nota1 y nota2''.
El DSL ofrece la forma \texttt{promedio} que hace exactamente eso:
recibe \queesesto{una lista de columnas}, genera la suma y la divide entre la
cantidad de columnas automáticamente.

\begin{lstlisting}[caption={Tabla con columna calculada.}]
(def-table tabla-notas
  ((nombre "") (nota1 "") (nota2 "") (promedio ""))
  :cell-height 1
  :cell-width  1
  :computed
    ((promedio (promedio (col nota1) (col nota2)))))
\end{lstlisting}

\notaparaelautor{Si intento ejecutar eso me da el mismo error de arriba... Que tiene sentido, porque no en lo que he hecho no está definido \texttt{def-table}.}

\noindent Cada pieza significa:

\begin{itemize}
  \item \texttt{(col nota1)} --- toma el valor de la columna
    \texttt{nota1} \comment{en la fila actual}{¿qué es la fila actual?}.
  \item \texttt{(promedio (col nota1) (col nota2))} --- suma nota1 y
    nota2, y divide entre 2 (porque recibe dos argumentos).
\end{itemize}

\begin{infobox}{¿Cómo sabe \texttt{promedio} entre cuánto dividir?}
  \texttt{promedio} cuenta automáticamente los argumentos que recibe.
  Con \texttt{(promedio (col n1) (col n2))} genera
  \texttt{((n1+n2)/2)}; si recibe tres columnas genera
  \texttt{((n1+n2+n3)/3)}.  No hace falta indicar el divisor.
\end{infobox}

\notaparaelautor{Ahora mismo te lo tengo que creer... Si tuviéramos la variante de crear el libro con esa hoja y esa tabla, pudiéramos comprobarlo \textit{de verdad}.}

\agregaesto{Este sería un buen momento para poner un cuadrito chulo de eso 8-) en el que comentes qué otras posibles funciones se pueden usar en lugar de promedio... Por ejemplo, si quiero la suma o si quiero comprobar si nota1 es mayor que nota2.}

\agregaesto{TRANSICIÓN A LO QUE SIGUE.}

\section{Resaltado condicional}

Ahora queremos que los nombres de las personas cuyo promedio supere el
mínimo aparezcan en rojo.  En lenguaje natural: ``si el promedio de
esta fila es mayor que nota-minima, colorea el nombre''.

La opción \texttt{:render} de \texttt{def-table} acepta una forma
\texttt{conditional-rendering} que especifica la condición y las
columnas a colorear.

\begin{lstlisting}[caption={Tabla con columna calculada y resaltado.}]
(def-table tabla-notas
  ((nombre "") (nota1 "") (nota2 "") (promedio ""))
  :cell-height 1
  :cell-width  1
  :computed
    ((promedio (promedio (col nota1) (col nota2))))
  :render
    (conditional-rendering
      :condition (gt (col promedio) (param nota-minima))
      :target-columns (nombre)))
\end{lstlisting}

\notaparaelautor{:-o... buenísimo esto :-o... lo entendí a la primera :-).}

\noindent Al leer la condición en voz alta:

\begin{itemize}
  \item \texttt{gt} significa ``mayor que'' (\emph{greater than}).
  \item \texttt{(col promedio)} es el promedio de la fila actual.
  \item \texttt{(param nota-minima)} es el valor fijo que definiremos.
  \item \texttt{:target-columns (nombre)} indica qué columna se colorea
    cuando la condición se cumple.
\end{itemize}

\notaparaelautor{Aquí vendría el mismo comentario que arriba, comenta por aquí (quizás en un cuadrito lindo) qué otras opciones existen para las comparaciones y cómo se pueden usar los operadores lógicos de conjunción, disyución y negación.}

\notaparaelautor{Y lo otro que sería bueno comentar aquí es qué hacer cuando quieres que se resalten más de una columna, o cuando quieres tener dos condicionales...}

\notaparaelautor{Bueno, en realidad, queda un problema colgáo que es el (param nota-minima) ese... del que no has dicho... la buena noticia es que eso se puede resolver cambiando el orden del tutorial...}

\agregaesto{TRANSICIÓN A LO QUE SIGUE.}

\section{El parámetro y la generación}

\comment{Con la tabla definida, solo falta ensamblar el libro}{esto sería bueno cambiarlo con la nueva idea, de hacerlo todo paso a paso... incluso, antes de los parámetros, puedes presentar cómo insertar solo un texto, y después dices que puede ser un parámero. }, proporcionar los
parámetros y \queesesto{lanzar la generación}.

\comment{Los parámetros}{¿los parámetros de qué?} se pasan \comment{al instanciar la tabla con \texttt{tabla}}{Esto debería estar explicado por allá arriba, con el ejemplito bobo...}
mediante la clave \texttt{:params}.  Cada parámetro es un par
\texttt{(nombre valor)}.

\begin{lstlisting}[caption={Libro completo con parámetros.}]
(libro tutorial-basico
  :filename "Tutorial_Basico.xlsx"
  :hojas (list
    (hoja "Notas"
      (tabla tabla-notas
        :data *DATOS*
        :params ((nota-minima 80))))))

(xl-generate tutorial-basico "tutorial-basico.py")
(xl-run-generated "tutorial-basico.py")
\end{lstlisting}

\notaparaelautor{Aquí me sigue pareciendo una mala idea que el libro tenga un \texttt{:filename}. Eso se debería cambiar... revisa los comentarios que te puse sobre eso en el otro documento.}

\begin{infobox}{El concepto de parámetro}
  Un \texttt{param} es un valor fijo con nombre que se escribe una sola
  vez y \comment{se referencia desde cualquier expresión}{¿qué \textit{scope} tienen esos parámetros? ¿son del libro, de la hoja, de la tabla?} con \texttt{(param
  nombre)}.  Es la manera que tiene el DSL de representar constantes
  que puedan cambiar entre ejecuciones sin tocar las fórmulas.  En este
  ejemplo \texttt{(param nota-minima)} vale 80.

  En la implementación concreta del generador a Excel, \comment{los
    parámetros se traducen a celdas fijas de la hoja (por ejemplo
    \texttt{\$F\$1})}{¿Y esto lo defino yo como usuario del DSL o
    alguien lo especificó por mí? Eso sería bueno decirlo.} para que las fórmulas puedan
  referenciarlos con referencias absolutas.  \queesesto{El backend de Python} los
  escribe directamente en \comment{esas celdas}{¿cuáles celdas?}, \comment{sin mezclarlos con las filas de
  datos de la tabla}{¿y si hay una tabla que tiene columnas que llegan hasta F?}.
\end{infobox}


\agregaesto{TRANSICIÓN A LO QUE SIGUE.}

\section{El programa completo}

A continuación se muestra el contenido íntegro del archivo
\texttt{tutorial-basico.lisp}:

\begin{lstlisting}[caption={tutorial-basico.lisp --- programa completo.}]
(load (merge-pathnames "dsl-directo.lisp" *load-truename*))

(defparameter *DATOS*
  '(("" "" "" "")
    ("" "" "" "")
    ("Nombre" "Nota 1" "Nota 2" "Promedio")
    ("Ana"     "85" "90" "")
    ("Luis"    "70" "75" "")
    ("Elena"   "92" "88" "")
    ("Carlos"  "60" "65" "")
    ("Sofia"   "78" "82" "")
    ("Pedro"   "95" "91" "")))

(def-table tabla-notas
  ((nombre "") (nota1 "") (nota2 "") (promedio ""))
  :cell-height 1
  :cell-width  1
  :computed
    ((promedio (promedio (col nota1) (col nota2))))
  :render
    (conditional-rendering
      :condition (gt (col promedio) (param nota-minima))
      :target-columns (nombre)))

(libro tutorial-basico
  :filename "Tutorial_Basico.xlsx"
  :hojas (list
    (hoja "Notas"
      (tabla tabla-notas
        :data *DATOS*
        :params ((nota-minima 80))))))

(xl-generate tutorial-basico "tutorial-basico.py")
(xl-run-generated "tutorial-basico.py")
\end{lstlisting}

\notaparaelautor{En el render tuve una duda... quise que me resaltara las columnas nombre y promedio... pero eso no pasó... lo que hice fue ponerle el nombre de la columna promedio, pero eso no funcionó... ¿tengo que poner dos formatos condicionales?}


\section{Ejecución}

Para ejecutar el tutorial:

\begin{lstlisting}[language=bash, numbers=none, backgroundcolor=\color{gray!10},
                   basicstyle=\ttfamily\small]
cd tutorial/
sbcl --script tutorial-basico.lisp
\end{lstlisting}

Esto cargará el DSL, construirá el AST, generará
\texttt{tutorial-basico.py} y lo ejecutará para crear
\texttt{Tutorial\_Basico.xlsx}.

\section{Resultado esperado}

El archivo \texttt{Tutorial\_Basico.xlsx} contiene una hoja llamada
``Notas'' con la siguiente estructura:

\begin{center}
\begin{tabular}{ll}
\toprule
\textbf{Rango} & \textbf{Contenido} \\
\midrule
Filas 3--9, columnas A--D & Tabla con 6 personas (fila 3: encabezados) \\
Columna D, filas 4--9     & Promedio calculado con fórmula \\
Columna A, filas 4--9     & Nombres en rojo si promedio $> 80$ \\
\bottomrule
\end{tabular}
\end{center}

Con los datos de ejemplo, los nombres de Ana, Elena, Sofía y Pedro
aparecen resaltados en rojo porque su promedio supera 80.  Luis y Carlos
quedan con formato normal.

\section{Resumen de formas utilizadas}

\begin{center}
\begin{tabular}{>{\ttfamily}p{5cm}p{8cm}}
\toprule
\textbf{Forma} & \textbf{Para qué sirve} \\
\midrule
(def-table nombre (cols...) opts...) & Define una plantilla de tabla \\
:computed ((col expr)) & Columna calculada con una expresión \\
(conditional-rendering :condition :target-columns) & Regla de resaltado condicional \\
(promedio col1 col2 ...) & Promedia las columnas indicadas \\
(gt a b) & Verdadero si a es mayor que b \\
(col nombre) & Valor de la columna en la fila actual \\
(param nombre) & Valor de un parámetro fijo \\
:params ((nombre valor)...) & Define parámetros al instanciar la tabla \\
(libro nombre ...) & Define el libro completo \\
(hoja nombre tablas ...) & Define una hoja \\
(tabla plantilla :data :params) & Instancia una tabla con datos y parámetros \\
(xl-generate wb archivo) & Genera el script Python \\
(xl-run-generated archivo) & Ejecuta el script y produce el .xlsx \\
\bottomrule
\end{tabular}
\end{center}

\end{document}
